\documentclass[12pt,a4paper]{article}

\usepackage[utf8]{inputenc}
\usepackage[brazil]{babel}
\usepackage[T1]{fontenc}
\usepackage[top=2.5cm,bottom=2.5cm,left=3cm,right=2cm]{geometry}
\usepackage{hyperref}
\usepackage{booktabs}
\usepackage{xcolor}
\usepackage{listings}
\usepackage{enumitem}

\hypersetup{
colorlinks=true,
linkcolor=blue,
filecolor=magenta,      
urlcolor=blue,
pdftitle={Documentação da Sala de Situação de Saúde - Backend},
}

% Configuração para trechos de código e JSON
\definecolor{codegreen}{rgb}{0,0.6,0}
\definecolor{codegray}{rgb}{0.5,0.5,0.5}
\definecolor{codepurple}{rgb}{0.58,0,0.82}
\definecolor{backcolour}{rgb}{0.97,0.97,0.97}

\lstset{
backgroundcolor=\color{backcolour},   
commentstyle=\color{codegreen},
keywordstyle=\color{blue},
numberstyle=\tiny\color{codegray},
stringstyle=\color{codepurple},
basicstyle=\ttfamily\footnotesize,
breakatwhitespace=false,         
breaklines=true,                 
keepspaces=true,                 
numbers=left,                    
numbersep=5pt,                  
showspaces=false,                
showstringspaces=false,
showtabs=false,                  
tabsize=2,
frame=single
}

\title{\textbf{Sala de Situação de Saúde (NSS) -- Backend}\\\large Documentação Técnica da API e Arquitetura}
\author{Núcleo de Sala de Situação de Saúde}
\date{\today}

\begin{document}

\maketitle

\begin{abstract}
Este documento descreve a arquitetura, pilha de tecnologias, estrutura organizacional de código, endpoints de API e instruções operacionais para o backend da Sala de Situação de Saúde (NSS), disponível no repositório \href{https://github.com/ArtursPereira/Site-Sala-de-Situa-o-de-Saude-Java}{Site-Sala-de-Situa-o-de-Saude-Java}.
\end{abstract}

\tableofcontents
\newpage

\section{Visão Geral e Status do Projeto}

O backend da plataforma de \textbf{Sala de Situação de Saúde (NSS)} é desenvolvido com Java 21 e o ecossistema Spring Boot 3. A aplicação tem como responsabilidade gerenciar a persistência de dados em banco relacional, garantir autenticação e autorização seguras e fornecer APIs REST para integração com os clientes web e pipelines de dados.

Atualmente, o projeto encontra-se em estágio inicial de desenvolvimento, com o módulo de \textbf{Autenticação e Gestão de Usuários} funcional:
\begin{itemize}[noitemsep]
\item Cadastro de novos usuários;
\item Autenticação via login com emissão de token JWT;
\item Hashing seguro de senhas via Argon2;
\item Consulta e listagem de usuários cadastrados.
\end{itemize}

\section{Pilha Tecnológica}

A arquitetura adota as seguintes tecnologias:
\begin{itemize}[noitemsep]
\item \textbf{Linguagem:} Java 21
\item \textbf{Framework:} Spring Boot 3.x (Spring Web, Spring Data JPA, Spring Security)
\item \textbf{Banco de Dados:} PostgreSQL
\item \textbf{Controle de Migrações:} Flyway
\item \textbf{Segurança:} JSON Web Tokens (JJWT 0.12.6) e Argon2 Password Encoder
\item \textbf{Mapeamento de Objetos:} MapStruct 1.6.3
\item \textbf{Produtividade:} Project Lombok
\item \textbf{Gerenciador de Dependências:} Apache Maven
\item \textbf{Contêineres:} Docker e Docker Compose
\end{itemize}

\section{Estrutura de Diretórios}

O projeto é estruturado em camadas sob o pacote base \texttt{com.example.demo}:

\begin{lstlisting}[language={}, numbers=none]
src/main/java/com/example/demo/
├── controller/     # Controladores REST (/auth, /NSS/users)
├── DTO/            # Objetos de transferência de dados (requests, responses, mappers)
├── entity/         # Entidades de persistência JPA (ex.: User)
├── repository/     # Interfaces de acesso a dados via Spring Data JPA
├── security/       # Filtros JWT, SecurityConfig e UserDetailsService
└── service/        # Regras de negócio e serviços da aplicação
\end{lstlisting}

\section{Modelo de Usuário e Autenticação}

A aplicação opera sob o modelo \textit{stateless}, exigindo o cabeçalho \texttt{Authorization: Bearer <token>} para endpoints protegidos. As senhas são processadas com o algoritmo \textbf{Argon2}.

A entidade principal \texttt{User} é composta pelos seguintes atributos:
\begin{itemize}[noitemsep]
\item \texttt{nome}: Nome completo do usuário;
\item \texttt{email}: E-mail de identificação corporativa ou pessoal (chave de login);
\item \texttt{password}: Hash da senha;
\item \texttt{matricula}: Registro de identificação institucional;
\item \texttt{cargo}: Perfil ou papel operacional desempenhado no núcleo.
\end{itemize}

\section{Endpoints da API REST}

\subsection{Autenticação (\texttt{/auth})}

\begin{table}[h]
\centering
\begin{tabular}{@{}llll@{}}
\toprule
\textbf{Método} & \textbf{Endpoint} & \textbf{Acesso} & \textbf{Descrição} \\ \midrule
POST & \texttt{/auth/register} & Público & Registro de novos usuários \\
POST & \texttt{/auth/login} & Público & Autenticação e emissão do token JWT \\ \bottomrule
\end{tabular}
\caption{Rotas de autenticação da aplicação.}
\end{table}

\noindent \textbf{Exemplo de Payload para Cadastro (\texttt{POST /auth/register}):}
\begin{lstlisting}
{
"nome": "Nome do Profissional",
"email": "usuario@saude.gov.br",
"password": "senhaSegura123",
"matricula": "123456",
"cargo": "ANALISTA"
}
\end{lstlisting}

\noindent \textbf{Exemplo de Payload para Login (\texttt{POST /auth/login}):}
\begin{lstlisting}
{
"email": "usuario@saude.gov.br",
"password": "senhaSegura123"
}
\end{lstlisting}

\noindent \textbf{Resposta esperada:}
\begin{lstlisting}
{
"token": "eyJhbGciOiJIUzI1NiIsInR5cCI6IkpXVCJ9..."
}
\end{lstlisting}

\subsection{Gestão de Usuários (\texttt{/NSS/users})}

\begin{table}[h]
\centering
\begin{tabular}{@{}llll@{}}
\toprule
\textbf{Método} & \textbf{Endpoint} & \textbf{Acesso} & \textbf{Descrição} \\ \midrule
GET & \texttt{/NSS/users} & Autenticado & Retorna todos os usuários cadastrados \\
GET & \texttt{/NSS/users/\{id\}} & Autenticado & Retorna os detalhes do usuário especificado \\ \bottomrule
\end{tabular}
\caption{Rotas de consulta e gestão de usuários.}
\end{table}

\section{Instalação e Execução}

\subsection{Pré-requisitos}
\begin{itemize}[noitemsep]
\item Java Development Kit (JDK) 21;
\item Docker e Docker Compose instalados;
\item Maven (opcional, o wrapper \texttt{mvnw} está incluso no repositório).
\end{itemize}

\subsection{Inicialização dos Serviços}

\noindent 1. Inicialize os contêineres de suporte (PostgreSQL):
\begin{lstlisting}[language=bash, numbers=none]
docker compose up -d
\end{lstlisting}

\noindent 2. Execute a aplicação via Spring Boot Wrapper:
\begin{lstlisting}[language=bash, numbers=none]
# Em ambientes Linux / macOS:
./mvnw clean spring-boot:run

# Em ambientes Windows:
mvnw.cmd clean spring-boot:run
\end{lstlisting}

Por padrão, a API estará acessível em \texttt{http://localhost:8080}.

\section{Roadmap e Integração com o Ecossistema NSS}

\begin{itemize}
\item \textbf{Controle de Acesso Baseado em Perfis (RBAC):} Formalização das \textit{roles} institucionais e restrições de rotas por perfil.
\item \textbf{Integração com o Frontend:} Conexão com o repositório \href{https://github.com/Zadoque/NSS-front-end}{NSS Front-end}.
\item \textbf{Módulos de Negócio em Saúde:} Implementação das entidades para vigilância em saúde, monitoramento de leitos e indicadores epidemiológicos.
\item \textbf{Infraestrutura e CI/CD:} Integração com o fluxo contínuo de implantação descrito em \href{https://github.com/Zadoque/NSS-deployment}{NSS Deployment}.
\end{itemize}

\end{document}